\chapter{Introducción a Haskell}

\section{¿Qué es Haskell?}
Haskell es un lenguaje funcional puro, de propósito general. Se utiliza en verificación de hardware, telecomunicaciones, comercio electrónico, etc.

\subsection{Características principales}
\begin{itemize}
  \item \textbf{Transparencia referencial}: el resultado depende solo de los argumentos.
  \item \textbf{Funciones como ciudadanos de primera clase}: se pasan, devuelven y almacenan.
  \item \textbf{Evaluación perezosa}: se evalúa solo cuando se necesita.
  \item \textbf{Tipos fuerte y estático}: verificación en compilación.
  \item \textbf{Inferencia de tipos}.
  \item \textbf{Polimorfismo paramétrico y sobrecarga}.
\end{itemize}

\section{Tipos Básicos}

\subsection{Enteros}
\begin{table}[h]
\centering
\begin{tabular}{|l|p{5cm}|p{5cm}|}
\hline
\textbf{Característica} & \textbf{Int} & \textbf{Integer} \\
\hline
Tamaño & Fijo & Arbitrario \\
Rango & Limitado & Ilimitado \\
Rendimiento & Rápido & Lento \\
Riesgo & Desbordamiento & Ninguno \\
\hline
\end{tabular}
\caption{Comparación entre Int e Integer}
\end{table}

\subsection{Reales}
\texttt{Float} (7 dígitos), \texttt{Double} (15-16 dígitos).

\subsection{Bool, Char, String}
\texttt{Bool}: \texttt{True}, \texttt{False}. \texttt{Char}: \texttt{'a'}. \texttt{String} es \texttt{[Char]}.

\section{Definición de funciones}

\begin{lstlisting}
areaRectangulo :: Int -> Int -> Int
areaRectangulo b h = b * h

maximo :: Int -> Int -> Int
maximo x y | x <= y = y
           | otherwise = x

factorial :: Integer -> Integer
factorial 0 = 1
factorial n = n * factorial (n - 1)
\end{lstlisting}

\section{Listas}

\begin{lstlisting}
[1,2,3] :: [Int]
[1..10] -- [1,2,3,4,5,6,7,8,9,10]
[1,3..10] -- [1,3,5,7,9]

dobles :: [Int] -> [Int]
dobles xs = [2*x | x <- xs]

map :: (a -> b) -> [a] -> [b]
map (+1) [1,2,3] -- [2,3,4]
\end{lstlisting}

\section{Funciones de Orden Superior}

\begin{lstlisting}
dosVeces :: (a -> a) -> a -> a
dosVeces f x = f (f x)

(.) :: (b -> c) -> (a -> b) -> (a -> c)
(g . f) x = g (f x)
\end{lstlisting}

\section{Tuplas}

\begin{lstlisting}
(1, "hola") :: (Int, String)
fst (1, "hola") -- 1
snd (1, "hola") -- "hola"
\end{lstlisting}

\section{Ejercicios}
(Ver el texto completo.)